Un thermomètre infrarouge mesure la puissance rayonnée par unité de surface d'un objet et en déduit la température en supposant une émissivité égale à 1. On vise une plaque métallique à \qty{\v{Tv}}{\degreeCelsius}, dont l'émissivité réelle est de \num{\v{em}}. On néglige le rayonnement ambiant réfléchi par la plaque.

\begin{enumerate}[label=\alph*)]
    \item Quelle puissance par mètre carré la plaque rayonne-t-elle ?
    \item Quelle température le thermomètre affiche-t-il, en degrés Celsius ?
    \item L'erreur est-elle une surestimation ou une sous-estimation ? Comment la corriger ?
\end{enumerate}
